\documentclass[aspectratio=169,10pt]{beamer}

\usetheme{Madrid}
\usecolortheme{default}
\setbeamertemplate{navigation symbols}{}
\setbeamertemplate{footline}[frame number]

\usepackage[T1]{fontenc}
\usepackage[utf8]{inputenc}
\usepackage{booktabs}
\usepackage{array}
\usepackage{graphicx}
\usepackage{hyperref}

\title[Thai Stock DRL Project]{Deep Reinforcement Learning for Thai Stock Portfolio Allocation}
\subtitle{Comprehensive 100-Slide Project Summary}
\author{Final Project Progress Report}
\institute{BistKA HPC Project: \texttt{/lustrefs/project/25sfcs03/drl\_thai\_stock}}
\date{June 9, 2026}

\begin{document}

% 1
\begin{frame}
  \titlepage
\end{frame}

% 2
\begin{frame}{How To Read This Deck}
\begin{itemize}
  \item This deck explains what has been built, tested, and learned.
  \item It uses simple language first, then adds technical details.
  \item The results are pilot results, not final market-performance claims.
  \item The deck is organized from motivation to data, model, results, and next steps.
\end{itemize}
\end{frame}

% 3
\begin{frame}{Project In One Sentence}
\begin{itemize}
  \item The project builds a reproducible research pipeline for testing reinforcement learning portfolio allocation on Thai stocks.
  \item It compares DRL agents with simple baselines.
  \item It adds multi-source features step by step.
  \item It runs and validates the work on BistKA HPC.
\end{itemize}
\end{frame}

% 4
\begin{frame}{Main Research Question}
\begin{itemize}
  \item Can a DRL portfolio agent improve long-term profit and risk control for Thai stocks?
  \item Can it beat simple baselines such as equal weight and buy-and-hold?
  \item Can extra information such as index, sector, macro, fundamentals, and sentiment help?
  \item Can the experiment be made reproducible enough to trust the result?
\end{itemize}
\end{frame}

% 5
\begin{frame}{Why This Project Matters}
\begin{itemize}
  \item Financial ML experiments can look good if data leakage or weak baselines are ignored.
  \item This project focuses on careful experimental wiring.
  \item It records data provenance and validation status.
  \item It shows what works now and what still needs better data.
\end{itemize}
\end{frame}

% 6
\begin{frame}{Final Status At A Glance}
\begin{itemize}
  \item Local repository scaffold is complete.
  \item HPC project directory and conda environment are complete.
  \item Latest verified HPC test suite passed with 53 tests.
  \item Final archive with ticker-universe CSV support has been created on HPC.
\end{itemize}
\end{frame}

% 7
\begin{frame}{Important Caveat}
\begin{itemize}
  \item The current real-data experiments use a five-stock starter universe.
  \item These results validate the pipeline, not the whole Thai market.
  \item Historical sentiment and official macro data do not yet overlap the 2021-2024 model window.
  \item Stronger performance claims require broader and licensed data.
\end{itemize}
\end{frame}

% 8
\begin{frame}{What Has Been Built}
\begin{itemize}
  \item Data ingestion for synthetic and real OHLCV data.
  \item Feature engineering for technical and contextual features.
  \item A Gymnasium-compatible trading environment.
  \item DRL training, baselines, evaluation, plots, and result aggregation.
  \item HPC job scripts and repeatable validation commands.
\end{itemize}
\end{frame}

% 9
\begin{frame}{Main Project Folder}
\begin{itemize}
  \item HPC working folder:
  \item \texttt{/lustrefs/project/25sfcs03/drl\_thai\_stock}
  \item Local working folder:
  \item \texttt{/Users/nathawat/Documents/final\_project\_plan}
  \item GitHub remote:
  \item \texttt{github.com/pppppppppppp-North/drl}
\end{itemize}
\end{frame}

% 10
\begin{frame}{Current Repository Shape}
\begin{itemize}
  \item \texttt{config/}: experiment settings.
  \item \texttt{src/}: ingestion, features, environment, training, experiments.
  \item \texttt{tests/}: unit and pipeline tests.
  \item \texttt{docs/}: methodology, data sources, and artifact notes.
  \item \texttt{slurm/}: HPC job templates.
\end{itemize}
\end{frame}

% 11
\begin{frame}{HPC Setup Completed}
\begin{itemize}
  \item BistKA project directory was created.
  \item Python environment was created at \texttt{/lustrefs/project/25sfcs03/envs/drl\_env}.
  \item Code, configs, docs, tests, and selected outputs were synced to HPC.
  \item File-separation snapshots track uploaded files versus generated files.
\end{itemize}
\end{frame}

% 12
\begin{frame}{Latest HPC Validation}
\begin{itemize}
  \item Latest full HPC test suite: 53 tests passed.
  \item Ticker-universe focused tests: 4 tests passed.
  \item Historical sentiment importer tests: 3 tests passed.
  \item Official macro CSV importer tests: 3 tests passed.
  \item This means recent code additions are covered by automated tests.
\end{itemize}
\end{frame}

% 13
\begin{frame}{Project Phases}
\begin{enumerate}
  \item Build a reproducible scaffold.
  \item Add synthetic pilot data.
  \item Add real OHLCV starter data.
  \item Build trading environment and baselines.
  \item Train and compare DRL agents.
  \item Add multi-source features and validation studies.
\end{enumerate}
\end{frame}

% 14
\begin{frame}{Synthetic Pilot Purpose}
\begin{itemize}
  \item Synthetic data was used first to test the code safely.
  \item It avoids waiting for real data before checking the pipeline.
  \item It verifies that training, evaluation, and plots can run end to end.
  \item It is not used for final market conclusions.
\end{itemize}
\end{frame}

% 15
\begin{frame}{Synthetic Data Generator}
\begin{itemize}
  \item Implemented deterministic synthetic price generation.
  \item Generates daily stock-like OHLCV rows.
  \item Includes placeholder market return, macro change, and sentiment score.
  \item Useful for tests and smoke runs.
\end{itemize}
\end{frame}

% 16
\begin{frame}{Real Starter Universe}
\begin{itemize}
  \item The real OHLCV starter universe has five Thai tickers.
  \item \texttt{ADVANC.BK}, \texttt{AOT.BK}, \texttt{CPALL.BK}, \texttt{KBANK.BK}, \texttt{PTT.BK}.
  \item Data source is Yahoo Finance through \texttt{yfinance}.
  \item Licensed SET or SETSMART data is still preferred for final publication.
\end{itemize}
\end{frame}

% 17
\begin{frame}{Real OHLCV Feature Table}
\begin{itemize}
  \item Feature file: \texttt{data/processed/features\_real\_ohlcv.parquet}.
  \item Rows: 4,840.
  \item Tickers: 5.
  \item Columns: 21.
  \item Date range: 2021-01-04 to 2024-12-30.
\end{itemize}
\end{frame}

% 18
\begin{frame}{Train Validation Test Split}
\begin{center}
\begin{tabular}{lrrr}
\toprule
Split & Date range & Rows & Dates \\
\midrule
Train & 2021-01-04 to 2023-05-29 & 2900 & 580 \\
Validation & 2023-05-30 to 2024-03-12 & 970 & 194 \\
Test & 2024-03-13 to 2024-12-30 & 970 & 194 \\
\bottomrule
\end{tabular}
\end{center}
\begin{itemize}
  \item The split is chronological.
  \item Future data is not mixed into earlier training periods.
\end{itemize}
\end{frame}

% 19
\begin{frame}{Data Quality Checks}
\begin{itemize}
  \item Data-quality reports were generated for the real feature tables.
  \item Ticker coverage and missingness are tracked.
  \item The real OHLCV starter table has complete five-ticker panel coverage.
  \item Leakage checks verify duplicate rows and split boundaries.
\end{itemize}
\end{frame}

% 20
\begin{frame}{Data Manifest}
\begin{itemize}
  \item \texttt{data/data\_manifest.csv} records source provenance.
  \item It records access method, frequency, date range, symbols, columns, and missing-value rate.
  \item It separates synthetic, Yahoo, BOT, and user-provided import paths.
  \item This improves reproducibility and honesty about source limits.
\end{itemize}
\end{frame}

% 21
\begin{frame}{Ticker Universe CSV Support}
\begin{itemize}
  \item Added \texttt{src/data/ticker\_universe.py}.
  \item Real ingesters can now read tickers from a CSV or text file.
  \item This avoids hard-coding all tickers in YAML.
  \item It prepares the project for a broader SET50 or licensed universe file.
\end{itemize}
\end{frame}

% 22
\begin{frame}{Starter Universe CSV}
\begin{itemize}
  \item Added \texttt{data/reference/thai\_starter\_universe.csv}.
  \item It contains the current five starter tickers.
  \item Added example config \texttt{config/real\_ohlcv\_universe\_csv.yaml}.
  \item The same mechanism can load a larger curated ticker list later.
\end{itemize}
\end{frame}

% 23
\begin{frame}{Technical Features}
\begin{itemize}
  \item 1-day, 5-day, and 20-day returns.
  \item 20-day rolling volatility.
  \item 10-day and 30-day moving-average ratios.
  \item RSI, MACD, Bollinger z-score, ATR, and volume ratio.
\end{itemize}
\end{frame}

% 24
\begin{frame}{Why Technical Features Are Used}
\begin{itemize}
  \item Raw prices alone are hard for a model to interpret.
  \item Returns describe recent price movement.
  \item Volatility describes risk.
  \item Momentum indicators describe trend and overbought or oversold behavior.
\end{itemize}
\end{frame}

% 25
\begin{frame}{SET Market Context}
\begin{itemize}
  \item Added SET index ingestion using \texttt{\^{}SET.BK}.
  \item Raw file: \texttt{data/raw/prices\_market\_indices.csv}.
  \item Merged file: \texttt{features\_real\_ohlcv\_market.parquet}.
  \item This gives each stock a market-wide context.
\end{itemize}
\end{frame}

% 26
\begin{frame}{SET Index Features}
\begin{itemize}
  \item \texttt{set\_return\_1d}.
  \item \texttt{set\_return\_5d}.
  \item \texttt{set\_return\_20d}.
  \item \texttt{set\_volatility\_20d}.
  \item \texttt{set\_ma\_ratio\_20}.
\end{itemize}
\end{frame}

% 27
\begin{frame}{Sector Context}
\begin{itemize}
  \item Added pilot sector mapping for the five starter tickers.
  \item Mapping file: \texttt{data/reference/sector\_mapping\_thai\_pilot.csv}.
  \item Features include peer return, relative return, peer count, and sector flags.
  \item This is a pilot fallback until better sector-index data is available.
\end{itemize}
\end{frame}

% 28
\begin{frame}{Sector Mapping}
\begin{center}
\begin{tabular}{ll}
\toprule
Ticker & Sector \\
\midrule
ADVANC.BK & ICT \\
AOT.BK & Transportation \\
CPALL.BK & Commerce \\
KBANK.BK & Banking \\
PTT.BK & Energy \\
\bottomrule
\end{tabular}
\end{center}
\end{frame}

% 29
\begin{frame}{Sector Index Probe}
\begin{itemize}
  \item Added Yahoo sector-index candidate probe.
  \item Tested 139 candidate symbols on HPC.
  \item Only banking symbol \texttt{\^{}BANK} was usable.
  \item This proved that a complete public Yahoo sector-index source was not available.
\end{itemize}
\end{frame}

% 30
\begin{frame}{Macro Proxy Features}
\begin{itemize}
  \item Added Yahoo Finance macro proxy ingestion.
  \item Series: USDTHB, Brent, WTI, gold, and US 10-year yield.
  \item Raw file: \texttt{data/raw/prices\_macro\_yahoo.csv}.
  \item These are proxies, not official Thai macro releases.
\end{itemize}
\end{frame}

% 31
\begin{frame}{Macro Proxy Merge}
\begin{itemize}
  \item Macro proxy values are merged by date.
  \item Merge uses forward fill from the latest known observation.
  \item This avoids using future observations.
  \item Merged feature table has 4,840 rows and 68 columns.
\end{itemize}
\end{frame}

% 32
\begin{frame}{Fundamental Features}
\begin{itemize}
  \item Added Yahoo annual and quarterly financial statement ingestion.
  \item Raw file has 8,701 statement rows.
  \item Coverage spans five tickers and 13 period ends.
  \item Features are merged only after a 60-day reporting lag.
\end{itemize}
\end{frame}

% 33
\begin{frame}{Why A Reporting Lag Matters}
\begin{itemize}
  \item A financial statement is not known on the exact period-end date.
  \item Using it too early would create lookahead bias.
  \item The project delays fundamentals by 60 days.
  \item This is a conservative rule for pilot experiments.
\end{itemize}
\end{frame}

% 34
\begin{frame}{Official Macro Source Probe}
\begin{itemize}
  \item Added Bank of Thailand source probe.
  \item Targeted BOTWEBSTAT Leading Economic Indicator table.
  \item Raw file: \texttt{data/raw/bot\_official\_macro.csv}.
  \item Release-date-aware feature merge is implemented.
\end{itemize}
\end{frame}

% 35
\begin{frame}{Official Macro Limitation}
\begin{itemize}
  \item BOT web table probe returned rows dated 2025-11-01 to 2026-04-01.
  \item The modeling window is 2021-2024.
  \item Therefore the probe has no historical overlap.
  \item No official macro performance claim is made yet.
\end{itemize}
\end{frame}

% 36
\begin{frame}{Official Macro CSV Importer}
\begin{itemize}
  \item Added \texttt{src/data/ingest\_official\_macro\_csv.py}.
  \item Supports long and wide CSV formats.
  \item Can assign release dates when a source file lacks them.
  \item This lets a downloaded historical BOT export enter the pipeline later.
\end{itemize}
\end{frame}

% 37
\begin{frame}{Sentiment Source Probe}
\begin{itemize}
  \item Added Yahoo latest-news probe.
  \item Raw file: \texttt{data/raw/news\_yahoo\_latest.csv}.
  \item It returned only 11 rows.
  \item Dates were 2025-06-24 to 2026-03-31, outside the model window.
\end{itemize}
\end{frame}

% 38
\begin{frame}{Sentiment Merge Scaffold}
\begin{itemize}
  \item Daily sentiment extraction is implemented.
  \item Sentiment is merged backward-only by ticker and date.
  \item Maximum news age is seven days.
  \item This prevents future news from leaking into earlier stock rows.
\end{itemize}
\end{frame}

% 39
\begin{frame}{Historical Sentiment CSV Importer}
\begin{itemize}
  \item Added \texttt{src/data/ingest\_sentiment\_csv.py}.
  \item It normalizes user-provided historical news or sentiment CSV files.
  \item It can produce event-level news and daily ticker sentiment.
  \item It waits for a valid historical source instead of fabricating data.
\end{itemize}
\end{frame}

% 40
\begin{frame}{No-Lookahead Principle}
\begin{itemize}
  \item Features must only use information known at or before the trading date.
  \item Market and macro proxies are backward or forward-filled from known history.
  \item Fundamentals use a reporting lag.
  \item News and official macro data use release dates.
\end{itemize}
\end{frame}

% 41
\begin{frame}{Trading Environment}
\begin{itemize}
  \item Implemented \texttt{ThaiStockTradingEnv}.
  \item It follows the Gymnasium API.
  \item The agent chooses long-only portfolio weights.
  \item The environment tracks portfolio value through time.
\end{itemize}
\end{frame}

% 42
\begin{frame}{Action Space}
\begin{itemize}
  \item The action is a vector of weights, one per stock.
  \item Weights are normalized into a long-only budget.
  \item Short selling is not used in this pilot.
  \item Cash weight is tracked separately.
\end{itemize}
\end{frame}

% 43
\begin{frame}{Observation Space}
\begin{itemize}
  \item The environment gives the model a lookback window.
  \item The window contains the selected feature columns.
  \item This creates a stable tensor shape for DRL algorithms.
  \item Tests verify reset and step observation shapes.
\end{itemize}
\end{frame}

% 44
\begin{frame}{Portfolio Accounting}
\begin{itemize}
  \item The environment applies transaction costs.
  \item It tracks net return, turnover, and drawdown.
  \item It records per-ticker action weights.
  \item These values are saved for evaluation and plots.
\end{itemize}
\end{frame}

% 45
\begin{frame}{Reward Function}
\begin{itemize}
  \item Reward combines return, turnover penalty, and drawdown penalty.
  \item Current real-data settings:
  \item \texttt{return\_weight = 1.0}
  \item \texttt{turnover\_penalty = 0.05}
  \item \texttt{drawdown\_penalty = 0.1}
\end{itemize}
\end{frame}

% 46
\begin{frame}{Why Penalize Turnover}
\begin{itemize}
  \item High turnover means frequent portfolio changes.
  \item Frequent trading causes transaction costs.
  \item A turnover penalty discourages noisy trading.
  \item It makes the reward more realistic.
\end{itemize}
\end{frame}

% 47
\begin{frame}{Why Penalize Drawdown}
\begin{itemize}
  \item Drawdown measures loss from a previous peak.
  \item Investors care about large losses, not only final return.
  \item A drawdown penalty encourages risk control.
  \item This is important for long-term portfolio optimization.
\end{itemize}
\end{frame}

% 48
\begin{frame}{Implemented DRL Algorithms}
\begin{itemize}
  \item PPO is the main algorithm.
  \item A2C is implemented and compared.
  \item Support paths exist for DDPG, TD3, and SAC.
  \item PPO and A2C were the main validated comparison models.
\end{itemize}
\end{frame}

% 49
\begin{frame}{Why PPO}
\begin{itemize}
  \item PPO is widely used for continuous-control RL problems.
  \item It is more stable than many older policy-gradient methods.
  \item It works well with vectorized environments.
  \item It is a reasonable first DRL baseline for portfolio allocation.
\end{itemize}
\end{frame}

% 50
\begin{frame}{Baselines Implemented}
\begin{itemize}
  \item Equal weight.
  \item Random allocation.
  \item Momentum.
  \item Buy-and-hold per ticker.
  \item Moving-average crossover.
  \item Mean-variance.
\end{itemize}
\end{frame}

% 51
\begin{frame}{Why Baselines Matter}
\begin{itemize}
  \item A DRL model is only useful if it beats simple alternatives.
  \item Buy-and-hold is a strong and honest benchmark.
  \item Equal weight is simple, transparent, and hard to overfit.
  \item Baselines keep the research conclusion grounded.
\end{itemize}
\end{frame}

% 52
\begin{frame}{Training Outputs}
\begin{itemize}
  \item Saved model files.
  \item Equity curves.
  \item Metric CSV files.
  \item Action-weight logs.
  \item TensorBoard logs and plot-ready outputs.
\end{itemize}
\end{frame}

% 53
\begin{frame}{Evaluation Metrics}
\begin{itemize}
  \item Cumulative return.
  \item Annualized return.
  \item Sharpe ratio.
  \item Maximum drawdown.
  \item Final portfolio value.
  \item Turnover and active trading steps.
\end{itemize}
\end{frame}

% 54
\begin{frame}{Plotting Utility}
\begin{itemize}
  \item Implemented result plotting utility.
  \item It generates equity-curve plots.
  \item It generates drawdown plots.
  \item It generates turnover and action-weight plots.
  \item These are useful for report-ready figures.
\end{itemize}
\end{frame}

% 55
\begin{frame}{Synthetic Comparison Result}
\begin{center}
\begin{tabular}{lrr}
\toprule
Policy & Cum. return & Sharpe \\
\midrule
buy\_hold\_KBANK.BK & 0.580 & 3.219 \\
ma\_crossover & 0.258 & 1.881 \\
equal\_weight & 0.153 & 1.398 \\
a2c & 0.167 & 1.373 \\
ppo & -0.013 & -0.108 \\
\bottomrule
\end{tabular}
\end{center}
\end{frame}

% 56
\begin{frame}{Synthetic Result Interpretation}
\begin{itemize}
  \item Synthetic data confirmed the whole pipeline runs.
  \item Some simple baselines beat PPO in the synthetic comparison.
  \item This is acceptable because synthetic data is only a test bed.
  \item The important result is that training, evaluation, and aggregation worked.
\end{itemize}
\end{frame}

% 57
\begin{frame}{Real OHLCV Main Comparison}
\begin{center}
\begin{tabular}{lrrr}
\toprule
Policy & Cum. return & Sharpe & Max DD \\
\midrule
buy\_hold\_PTT.BK & 0.044 & 0.405 & -0.103 \\
buy\_hold\_ADVANC.BK & 0.010 & 0.097 & -0.099 \\
equal\_weight & -0.031 & -0.417 & -0.103 \\
a2c & -0.038 & -0.434 & -0.120 \\
ppo & -0.036 & -0.440 & -0.066 \\
\bottomrule
\end{tabular}
\end{center}
\end{frame}

% 58
\begin{frame}{Real OHLCV Interpretation}
\begin{itemize}
  \item Best validation policy was buy-and-hold PTT.
  \item PPO and A2C did not beat the strongest simple baseline.
  \item PPO had smaller max drawdown than A2C and equal weight.
  \item This is a useful and honest pilot outcome.
\end{itemize}
\end{frame}

% 59
\begin{frame}{Invalidated Run}
\begin{itemize}
  \item Run \texttt{real\_ohlcv\_compare\_2854} was invalidated.
  \item Reason: an old training entrypoint regenerated synthetic data into real paths.
  \item Corrected valid run is \texttt{real\_ohlcv\_compare\_2855}.
  \item This was recorded clearly in the checklist and manifest.
\end{itemize}
\end{frame}

% 60
\begin{frame}{Hyperparameter Search}
\begin{itemize}
  \item Added Optuna-style trial runner and SLURM array workflow.
  \item Search space includes learning rate, gamma, GAE lambda, entropy coefficient, rollout length, clip range, and batch size.
  \item Fixed search run: \texttt{optuna\_search\_fixed\_2863}.
  \item Earlier repeated-parameter search was invalidated.
\end{itemize}
\end{frame}

% 61
\begin{frame}{Best Short Optuna Trial}
\begin{center}
\begin{tabular}{lrrrr}
\toprule
Trial & Algo & Sharpe & Cum. return & Max DD \\
\midrule
4 & PPO & 0.280 & 0.020 & -0.069 \\
0 & PPO & -0.306 & -0.003 & -0.009 \\
1 & A2C & -0.520 & -0.043 & -0.127 \\
2 & PPO & -1.080 & -0.089 & -0.115 \\
3 & A2C & -1.708 & -0.003 & -0.004 \\
\bottomrule
\end{tabular}
\end{center}
\end{frame}

% 62
\begin{frame}{Longer Tuned PPO Run}
\begin{itemize}
  \item Best short PPO trial was promoted into a longer config.
  \item Longer run used 20,000 timesteps.
  \item Result: cumulative return -0.030 and Sharpe -0.403.
  \item The tuned parameters were not robust enough yet.
\end{itemize}
\end{frame}

% 63
\begin{frame}{Walk-Forward Validation}
\begin{itemize}
  \item Added rolling-window split generation.
  \item Starter settings: 504 training dates, 126 validation dates, 126 test dates.
  \item Step size is 126 dates.
  \item First pilot ran two short PPO windows on HPC.
\end{itemize}
\end{frame}

% 64
\begin{frame}{Walk-Forward Pilot Results}
\begin{center}
\begin{tabular}{lrr}
\toprule
Policy & Mean cum. return & Mean Sharpe \\
\midrule
buy\_hold\_ADVANC.BK & 0.084 & 1.468 \\
buy\_hold\_KBANK.BK & 0.064 & 1.157 \\
buy\_hold\_PTT.BK & 0.041 & 0.672 \\
equal\_weight & 0.021 & 0.633 \\
ppo & 0.004 & 0.058 \\
\bottomrule
\end{tabular}
\end{center}
\end{frame}

% 65
\begin{frame}{Walk-Forward Interpretation}
\begin{itemize}
  \item Walk-forward machinery is verified.
  \item The first PPO pilot produced small positive average test return.
  \item Simple buy-and-hold policies still performed better.
  \item More windows and richer data are needed for stronger conclusions.
\end{itemize}
\end{frame}

% 66
\begin{frame}{Ablation Study Purpose}
\begin{itemize}
  \item Ablations test which feature groups help.
  \item They prevent guessing about feature value.
  \item The pilot compared returns-only, technical, and technical-plus-context configs.
  \item It used short 1,000-timestep PPO jobs.
\end{itemize}
\end{frame}

% 67
\begin{frame}{PPO Ablation Summary}
\begin{center}
\begin{tabular}{lrr}
\toprule
Feature group & Cum. return & Sharpe \\
\midrule
Returns only & -0.010 & -0.453 \\
Technical + context & -0.056 & -0.824 \\
Technical only & -0.079 & -1.611 \\
\bottomrule
\end{tabular}
\end{center}
\end{frame}

% 68
\begin{frame}{Ablation Interpretation}
\begin{itemize}
  \item Ablation machinery is verified.
  \item Short PPO ablations did not beat buy-and-hold PTT.
  \item Returns-only PPO was least negative among the three PPO ablations.
  \item Stronger claims need broader data and longer validation.
\end{itemize}
\end{frame}

% 69
\begin{frame}{Regime And Stress Testing}
\begin{itemize}
  \item Added regime/stress-test runner.
  \item It evaluates baselines plus saved PPO/A2C models.
  \item It uses named market slices and high-volatility windows.
  \item Output includes slices, metrics, equity curves, and summary report.
\end{itemize}
\end{frame}

% 70
\begin{frame}{Regime Test Caveat}
\begin{itemize}
  \item Acute 2020 COVID crash cannot be evaluated.
  \item Reason: current real OHLCV starts on 2021-01-04.
  \item The available COVID-related slice is COVID recovery in 2021.
  \item This limitation is explicitly recorded.
\end{itemize}
\end{frame}

% 71
\begin{frame}{Best Policy By Regime}
\begin{center}
\begin{tabular}{llr}
\toprule
Regime & Best policy & Sharpe \\
\midrule
COVID recovery 2021 & ADVANC buy-hold & 2.376 \\
High volatility 1 & ADVANC buy-hold & -0.735 \\
High volatility 2 & A2C & 7.221 \\
Inflation hike 2022 & AOT buy-hold & 1.367 \\
Validation period & PTT buy-hold & 0.405 \\
\bottomrule
\end{tabular}
\end{center}
\end{frame}

% 72
\begin{frame}{Market-Context Pilot}
\begin{itemize}
  \item First short SET market-context validation run completed.
  \item Output directory: \texttt{real\_ohlcv\_market\_pilot\_20260531}.
  \item It verifies SET-index feature ingestion and merge.
  \item It still does not solve sector-relative data.
\end{itemize}
\end{frame}

% 73
\begin{frame}{Sector-Context Pilot}
\begin{itemize}
  \item First short sector-context validation run completed.
  \item Output directory: \texttt{real\_ohlcv\_sector\_pilot\_20260531}.
  \item PPO improved over the SET-market-context PPO pilot.
  \item It still did not beat the PTT buy-and-hold baseline.
\end{itemize}
\end{frame}

% 74
\begin{frame}{Macro-Proxy Pilot}
\begin{itemize}
  \item First short macro-proxy validation run completed.
  \item Output directory: \texttt{real\_ohlcv\_macro\_pilot\_20260531}.
  \item Macro-proxy PPO improved over sector-context PPO.
  \item It beat equal weight but not the strongest buy-and-hold baselines.
\end{itemize}
\end{frame}

% 75
\begin{frame}{Fundamentals Pilot}
\begin{itemize}
  \item First short fundamentals validation run completed.
  \item Output directory: \texttt{real\_ohlcv\_fundamentals\_pilot\_20260531}.
  \item PPO had lower drawdown than equal weight.
  \item But it still did not beat PTT or ADVANC buy-and-hold.
\end{itemize}
\end{frame}

% 76
\begin{frame}{Sentiment Scaffold Status}
\begin{itemize}
  \item News probe, daily scoring, and merge paths are implemented.
  \item Historical sentiment CSV importer is implemented and tested.
  \item Current Yahoo latest news has no 2021-2024 overlap.
  \item A valid historical news source is still needed.
\end{itemize}
\end{frame}

% 77
\begin{frame}{Official Macro Scaffold Status}
\begin{itemize}
  \item BOT source probe and release-date merge are implemented.
  \item Historical official macro CSV importer is implemented and tested.
  \item Current BOT web rows have no 2021-2024 overlap.
  \item A historical official macro export is still needed.
\end{itemize}
\end{frame}

% 78
\begin{frame}{Testing Strategy}
\begin{itemize}
  \item Unit tests cover environment behavior and feature logic.
  \item Ingestion tests cover source normalization.
  \item Experiment tests cover aggregation and config logic.
  \item HPC full-suite validation is used because local Python lacks some dependencies.
\end{itemize}
\end{frame}

% 79
\begin{frame}{Latest Test Progress}
\begin{center}
\begin{tabular}{lr}
\toprule
Milestone & Tests passed \\
\midrule
Sector-index probe & 43 \\
Official macro CSV importer & 46 \\
Historical sentiment CSV importer & 49 \\
Ticker-universe CSV support & 53 \\
\bottomrule
\end{tabular}
\end{center}
\end{frame}

% 80
\begin{frame}{File Separation On HPC}
\begin{itemize}
  \item User asked to separate uploaded files from files already on HPC.
  \item Added \texttt{scripts/file\_separation\_snapshot.py}.
  \item Latest snapshot: \texttt{\_file\_separation/20260531\_231548}.
  \item This helps distinguish code uploads from generated outputs.
\end{itemize}
\end{frame}

% 81
\begin{frame}{Final Archive}
\begin{itemize}
  \item Latest archive:
  \item \texttt{final\_archive/20260531\_final\_project\_pilot\_with\_ticker\_universe\_csv.tar.gz}.
  \item It contains source, configs, tests, docs, selected data, reports, and selected result outputs.
  \item It intentionally excludes large training caches and TensorBoard event files.
\end{itemize}
\end{frame}

% 82
\begin{frame}{GitHub Progress}
\begin{itemize}
  \item Project was pushed to GitHub.
  \item Recent commits include official macro CSV import, historical sentiment CSV import, and ticker-universe CSV support.
  \item Latest pushed commit before this slide deck work: \texttt{56a5a52}.
  \item Remote branch: \texttt{main}.
\end{itemize}
\end{frame}

% 83
\begin{frame}{Main Engineering Lesson}
\begin{itemize}
  \item The pipeline matters as much as the model.
  \item Without clean data splits, leakage checks, and baselines, DRL results are easy to overstate.
  \item This project now has a reliable experimental skeleton.
  \item The next bottleneck is data coverage, not only model code.
\end{itemize}
\end{frame}

% 84
\begin{frame}{Main Modeling Lesson}
\begin{itemize}
  \item PPO and A2C did not beat the strongest simple baselines in the starter experiment.
  \item This does not mean DRL is useless.
  \item It means the current data and feature set are not enough to justify a strong DRL claim.
  \item The result is still useful because it sets a realistic benchmark.
\end{itemize}
\end{frame}

% 85
\begin{frame}{What Worked Well}
\begin{itemize}
  \item End-to-end training and evaluation works.
  \item Baselines are implemented and comparable.
  \item Multi-source feature pipelines are modular.
  \item HPC testing and archiving workflow is repeatable.
  \item Limitations are documented instead of hidden.
\end{itemize}
\end{frame}

% 86
\begin{frame}{What Is Still Missing}
\begin{itemize}
  \item Broader curated or licensed stock universe.
  \item Complete sector-index or constituent source.
  \item Historical official macro source with release dates.
  \item Historical Thai news or sentiment source.
  \item Licensed fundamentals with stronger coverage.
\end{itemize}
\end{frame}

% 87
\begin{frame}{Why Results Are Called Pilot Results}
\begin{itemize}
  \item Only five real tickers are active.
  \item Some feature groups are scaffolds because historical source data is missing.
  \item Short 1,000-timestep pilots verify wiring, not final performance.
  \item The project is ready for stronger experiments once data improves.
\end{itemize}
\end{frame}

% 88
\begin{frame}{Data Sources Completed}
\begin{itemize}
  \item Synthetic pilot prices.
  \item Yahoo real OHLCV starter data.
  \item Yahoo SET index data.
  \item Yahoo macro proxy data.
  \item Yahoo statement fundamentals.
  \item Yahoo latest-news probe.
  \item BOT official macro probe.
\end{itemize}
\end{frame}

% 89
\begin{frame}{Import Paths Completed}
\begin{itemize}
  \item OHLCV ingestion.
  \item Market index ingestion.
  \item Macro proxy ingestion.
  \item Fundamentals ingestion.
  \item Latest-news probe.
  \item Historical sentiment CSV import.
  \item Historical official macro CSV import.
  \item Ticker-universe CSV import.
\end{itemize}
\end{frame}

% 90
\begin{frame}{Experiment Paths Completed}
\begin{itemize}
  \item Single training run.
  \item Algorithm comparison.
  \item Baseline evaluation.
  \item Hyperparameter search.
  \item Walk-forward pilot.
  \item Ablation pilot.
  \item Regime and stress tests.
\end{itemize}
\end{frame}

% 91
\begin{frame}{Report Artifacts Completed}
\begin{itemize}
  \item Final methodology/results markdown report.
  \item Final methodology/results LaTeX report.
  \item Data source documentation.
  \item Final artifact manifest.
  \item Progress checklist.
  \item Data-quality reports and selected figures.
\end{itemize}
\end{frame}

% 92
\begin{frame}{How To Explain The Result Simply}
\begin{itemize}
  \item The system works.
  \item The model is not yet better than simple buy-and-hold.
  \item The experiment is honest about this.
  \item The next improvement should focus on better and broader data.
\end{itemize}
\end{frame}

% 93
\begin{frame}{Recommended Next Step 1}
\begin{itemize}
  \item Replace the five-ticker starter universe with a broader curated list.
  \item Use the new ticker-universe CSV support.
  \item Rebuild OHLCV and feature tables.
  \item Re-run data-quality and leakage checks before training.
\end{itemize}
\end{frame}

% 94
\begin{frame}{Recommended Next Step 2}
\begin{itemize}
  \item Replace pilot sector mapping with a complete sector or constituent source.
  \item Then rebuild sector-relative features.
  \item This will make sector-relative signals more meaningful.
  \item It will also support broader universe experiments.
\end{itemize}
\end{frame}

% 95
\begin{frame}{Recommended Next Step 3}
\begin{itemize}
  \item Obtain historical official macro data with release dates.
  \item Use the official macro CSV importer.
  \item Rebuild official macro feature tables.
  \item Only then run official macro ablations.
\end{itemize}
\end{frame}

% 96
\begin{frame}{Recommended Next Step 4}
\begin{itemize}
  \item Obtain historical Thai news, disclosures, or sentiment scores.
  \item Use the historical sentiment CSV importer.
  \item Build daily ticker-level sentiment.
  \item Re-run sentiment feature merge and ablations.
\end{itemize}
\end{frame}

% 97
\begin{frame}{Recommended Next Step 5}
\begin{itemize}
  \item Expand walk-forward validation after better data is connected.
  \item Use more windows and longer training jobs.
  \item Compare PPO/A2C with every baseline in each window.
  \item Report average and per-window stability.
\end{itemize}
\end{frame}

% 98
\begin{frame}{Final Takeaway}
\begin{itemize}
  \item The project now has a working, tested, reproducible pilot pipeline.
  \item It includes data ingestion, features, DRL training, baselines, and evaluation.
  \item Current DRL results do not beat the strongest simple baseline.
  \item This is a strong foundation for the next data-driven experiment.
\end{itemize}
\end{frame}

% 99
\begin{frame}{What To Say In A Presentation}
\begin{itemize}
  \item "We built the full experimental system first."
  \item "We tested it on synthetic data and real Thai starter data."
  \item "We compared DRL with simple but strong baselines."
  \item "The honest result is that better data is needed before claiming DRL superiority."
\end{itemize}
\end{frame}

% 100
\begin{frame}{End}
\begin{itemize}
  \item Thank you.
  \item Full project code, configs, tests, reports, and archive paths are documented.
  \item Latest verified HPC test count: 53 passed.
  \item Latest final archive includes ticker-universe CSV support.
\end{itemize}
\end{frame}

\end{document}
